%--------------------------------------
% Streszczenie po polsku
%--------------------------------------
\cleardoublepage % Zaczynamy od nieparzystej strony
\streszczenie

Mutacje aminokwasowe w~białkach wirusowych stanowią kluczowy czynnik warunkujący zdolność patogenów do unikania odpowiedzi immunologicznej gospodarza. Zdolność do przewidywania przyszłych mutacji w~regionach epitopowych białek powierzchniowych wirusów ma istotne znaczenie dla projektowania szczepionek i~monitorowania epidemiologicznego. Niniejsza praca przedstawia konfigurowalny framework do predykcji mutacji aminokwasowych w~białkach wirusowych z~wykorzystaniem rekurencyjnych sieci neuronowych.

Zaproponowany potok przetwarzania obejmuje następujące etapy: parsowanie sekwencji w~formacie FASTA, kodowanie aminokwasów metodą ProtVec, klasteryzację sekwencji algorytmem K-means w~ramach poszczególnych okresów czasowych, łączenie klastrów między kolejnymi okresami na podstawie odległości centroidów oraz generowanie próbek treningowych z~wykorzystaniem okna przesuwnego. Framework udostępnia trzy architektury modeli predykcyjnych: standardową sieć LSTM, LSTM z~mechanizmem atencji temporalnej (Attention LSTM) oraz LSTM z~podwójnym mechanizmem atencji operującym zarówno w~wymiarze czasowym, jak i~wymiarze cech (Dual Attention LSTM).

Ewaluację przeprowadzono na dwóch zbiorach danych: sekwencjach białka S (Spike) wirusa SARS-CoV-2 pozyskanych z~bazy GISAID oraz sekwencjach białka hemagglutyniny (HA) wirusa Influenza A podtypu H1N1 z~bazy NCBI.
% TODO: Uzupełnić wyniki --- np.: Model Dual Attention LSTM osiągnął najwyższe
% wartości współczynnika korelacji Matthewsa (MCC) wynoszące odpowiednio
% X,XX dla SARS-CoV-2 i Y,YY dla Influenza, przewyższając zarówno model bazowy
% (regresja logistyczna), jak i standardową architekturę LSTM.
Zaprojektowany framework charakteryzuje się konfigurowalnością opartą na plikach YAML, reprodukowalnością eksperymentów oraz rozszerzalnością umożliwiającą zastosowanie do nowych białek wirusowych bez modyfikacji kodu źródłowego.

\slowakluczowe predykcja mutacji, SARS-CoV-2, Influenza, LSTM, mechanizm atencji, ProtVec, klasteryzacja
